\documentclass[
	% -- opções da classe memoir --
	article,	% indica que é um artigo acadêmico
	11pt,		% tamanho da fonte
	oneside,	% para impressão apenas no recto. Oposto a twoside
	a4paper,	% tamanho do papel. 
	% -- opções da classe abntex2 --
	%chapter=TITLE,   % títulos de capítulos convertidos em letras maiúsculas
	%section=TITLE,   % títulos de seções convertidos em letras maiúsculas
	%subsection=TITLE,  % títulos de subseções convertidos em letras maiúsculas
	%subsubsection=TITLE % títulos de subsubseções convertidos em letras maiúsculas
	% -- opções do pacote babel --
	english,	% idioma adicional para hifenização
	brazil,	% o último idioma é o principal do documento
	sumario=tradicional
	]{abntex2}

% ---
% PACOTES
% ---

% ---
% Pacotes fundamentais 
% ---
\usepackage{lmodern}														% Usa a fonte Latin Modern
\usepackage[T1]{fontenc}													% Selecao de codigos de fonte.
\usepackage[utf8]{inputenc}													% Codificacao do documento (conversão automática dos acentos)
\usepackage{indentfirst}													% Indenta o primeiro parágrafo de cada seção.
\usepackage{nomencl}														% Lista de simbolos
\usepackage{color}															% Controle das cores
\usepackage{graphicx}														% Inclusão de gráficos
\usepackage{microtype}														% para melhorias de justificação
% ---
\usepackage{titlesec}%[newparttoc]% Pacotes adicionais, usados apenas no âmbito do Modelo Canônico do abnteX2
\usepackage{lipsum}       % para geração de dummy text

\usepackage{tabularx}
\usepackage{array,makecell}
\usepackage{multirow}														% tabelas complexas
\usepackage{mdframed}														% ambiente de seções "boxed"
% ---
% ---
% ---
% Pacotes de citações
% ---
\usepackage[brazilian,hyperpageref]{backref}								% Paginas com as citações na bibl
\usepackage[alf]{abntex2cite}												% Citações padrão ABNT

% \newcounter{mydef}
% \renewcommand\themydef{\arabic{mydef} -- }
% \setcounter{mydef}{2}

% \mdtheorem[style=theoremstyle,theoremseparator= -- ]{definition}{#1 -- }
\newcounter{theo}
\setcounter{theo}{2}
\newenvironment{theo}[1][]{
\stepcounter{theo}
\ifstrempty{#1}
{\mdfsetup{
frametitle={
{Theorem~\thetheo};}}
}
{\mdfsetup{
frametitle={\thetheo~--~#1}}
}
\mdfsetup{frametitlerule=true}
% linewidth=2pt,topline=true,
% frametitleaboveskip=\dimexpr−\ht\strutbox\relax,}
 \begin{mdframed}[]\relax
 }{\end{mdframed}}

% \usepackage{siunitx}
% \newtheorem{theorem}{Theorem}[section]
% \renewcommand{\thetheorem}{\thesection--\arabic{theorem}}
% % \renewcommand{\thetheorem}{\thesection--\arabic{theorem}}
% \mdtheorem[style=theoremstyle]{definition}{Definition}[section]
% \mdtheorem[style=theoremstyle]{theorem}[definition]{Theorem}
\setheadfoot{3cm}{1cm}

\makepagestyle{meuestilo}
	%%cabeçalhos
	\makeevenhead{meuestilo} %%pagina par
		 {}
		 {\includegraphics[width=\textwidth]{fige.pdf}}
	{}
	\makeoddhead{meuestilo} %%pagina ímpar ou com oneside
		 {}
		 {\includegraphics[width=.9\textwidth]{fige.pdf}}
		 {}
% ---
% Configurações do pacote backref
% Usado sem a opção hyperpageref de backref
\renewcommand{\backrefpagesname}{Citado na(s) página(s):~}
% Texto padrão antes do número das páginas
\renewcommand{\backref}{}
% Define os textos da citação
\renewcommand*{\backrefalt}[4]{
	\ifcase #1 %
		Nenhuma citação no texto.%
	\or
		Citado na página #2.%
	\else
		Citado #1 vezes nas páginas #2.%
	\fi}%
% ---

% --- Informações de dados para CAPA e FOLHA DE ROSTO ---
\titulo{Título PIBEX 2023}
% \tituloestrangeiro{Canonical article template in \abnTeX: optional foreign title}
\local{Campus Ipojuca}

% \autor{
% Equipe \abnTeX\thanks{``Recomenda-se que os dados de vinculação e
% endereço constem em nota, com sistema de chamada próprio, diferente do sistema 
% adotado para citações no texto.'' \url{http://www.abntex.net.br/}} 
% \\[0.5cm] 
% Lauro César Araujo\thanks{``Constar currículo sucinto de cada autor, com 
% vinculação corporativa e endereço de contato.''} }

\data{Ipojuca, \today}
% ---
% ---
% Configurações de aparência do PDF final

% alterando o aspecto da cor azul
\definecolor{blue}{RGB}{41,5,195}

% informações do PDF
\makeatletter
\hypersetup{
			%pagebackref=true,
		pdftitle={\@title}, 
		pdfauthor={\@author},
			pdfsubject={Modelo de artigo científico com abnTeX2},
			pdfcreator={LaTeX with abnTeX2},
		pdfkeywords={abnt}{latex}{abntex}{abntex2}{atigo científico}, 
		colorlinks=true,          % false: boxed links; true: colored links
			linkcolor=blue,           % color of internal links
			citecolor=blue,           % color of links to bibliography
			filecolor=magenta,          % color of file links
		urlcolor=blue,
		bookmarksdepth=4
}
\makeatother
% --- 

% ---
% compila o indice
% ---
\makeindex
% ---

% ---
% Altera as margens padrões
% ---
\setlrmarginsandblock{2cm}{2cm}{*}
\setulmarginsandblock{4.5cm}{2cm}{*}
\checkandfixthelayout
% ---

% --- 
% Espaçamentos entre linhas e parágrafos 
% --- 

% O tamanho do parágrafo é dado por:
\setlength{\parindent}{0.0cm}

% Controle do espaçamento entre um parágrafo e outro:
\setlength{\parskip}{0.2cm}  % tente também \onelineskip

% Espaçamento simples
\SingleSpacing


% ----
% Início do documento
% ----
% \newcolumntype{s}{>{\hsize=.5\hsize}X}

\newcolumntype{a}{>{\hsize=.21\hsize}X}
\newcolumntype{e}{>{\hsize=1.050\hsize}X}
\newcolumntype{i}{>{\hsize=.35\hsize}X}
\newcolumntype{o}{>{\hsize=1.890\hsize}X}
\newcolumntype{u}{>{\hsize=2.240\hsize}X}
\newcolumntype{g}{>{\hsize=1.75\hsize}X}
\newcolumntype{s}{>{\hsize=0.5\hsize}X}

\begin{document}
\pagestyle{meuestilo}
\checkandfixthelayout
% Seleciona o idioma do documento (conforme pacotes do babel)
%\selectlanguage{english}
\selectlanguage{brazil}

% Retira espaço extra obsoleto entre as frases.
\frenchspacing 

% ----------------------------------------------------------
% ELEMENTOS PRÉ-TEXTUAIS
% ----------------------------------------------------------

%---
%
% Se desejar escrever o artigo em duas colunas, descomente a linha abaixo
% e a linha com o texto ``FIM DE ARTIGO EM DUAS COLUNAS''.
% \twocolumn[       % INICIO DE ARTIGO EM DUAS COLUNAS
%
%---

% página de titulo principal (obrigatório)
%\maketitle




% titulo em outro idioma (opcional)



% ----------------------------------------------------------
% ELEMENTOS TEXTUAIS
% ----------------------------------------------------------
\textual
\pagestyle{meuestilo}

\phantom{This text will be invisible}
\vfill

\begin{center}
\imprimirtitulo
\end{center}

\vfill

\begin{center}
Local de execução: \imprimirlocal
\end{center}

\vfill

\begin{center}
\imprimirdata
\end{center}
\newpage

\begin{center}
\begin{table}[h]
%\begin{tabular}{p{.3cm}p{1.9cm}p{1.1cm}p{3.9cm}p{0.3cm}p{1.1cm}p{5.1cm}}
\begin{tabularx}{\linewidth}{aeioaeu}

 &  & & & & &  \\\hline
\multicolumn{7}{|l|}{\textbf{1 -- DADOS DO PROJETO}}                                                                                             \\ \hline
\multicolumn{7}{|l|}{Título: \imprimirtitulo}                                \\ \hline
\multicolumn{7}{|l|}{Área Temática Principal}                                                                                                   \\ \hline
%\multicolumn{7}{|l|}{\textit{(assinalar apenas uma Área   Temática, aquela que melhor define o problema que deu origem ao projeto)}} \\ \hline
\multicolumn{1}{|c|}{}        & \multicolumn{3}{l|}{Comunicação}    & \multicolumn{1}{c|}{}          & \multicolumn{2}{l|}{Meio Ambiente}     \\ \hline
\multicolumn{1}{|c|}{}         & \multicolumn{3}{l|}{Cultura}        & \multicolumn{1}{c|}{}           & \multicolumn{2}{l|}{Saúde}             \\ \hline
\multicolumn{1}{|c|}{} &
	\multicolumn{3}{l|}{Direitos Humanos e Justiça} &
	\multicolumn{1}{c|}{X} &
	\multicolumn{2}{l|}{Tecnologia e Produção} \\ \hline
\multicolumn{1}{|c|}{}         & \multicolumn{3}{l|}{Educação}       & \multicolumn{1}{c|}{}           & \multicolumn{2}{l|}{Trabalho}          \\ \hline
\multicolumn{3}{|l|}{\begin{tabular}[c]{@{}l@{}}Local de realização/\\ Município\end{tabular}} &
	\multicolumn{4}{l|}{Microregião de SUAPE no Estado de Pernambuco, Brasil} \\ \hline
\multicolumn{3}{|l|}{Comunidade}                       & \multicolumn{4}{l|}{Comunidades dos municípios de Ipojuca e Cabo de Santo Agostinho} \\ \hline
\multicolumn{2}{|l|}{\begin{tabular}[c]{@{}l@{}}Beneficiários \\ Diretos\end{tabular}} &
	\multicolumn{2}{l|}{\begin{tabular}[c]{@{}l@{}}Moradores da \\ região de X\end{tabular}} &
	\multicolumn{2}{l|}{\begin{tabular}[c]{@{}l@{}}Beneficiários \\ Indiretos\end{tabular}} &
	\multicolumn{1}{l|}{\begin{tabular}[c]{@{}l@{}}Trabalhadores da \\ empresa X\end{tabular}} \\ \hline
\end{tabularx}
\end{table}
\end{center}

\begin{table}[h]
		\centering
		\begin{tabularx}{\textwidth}{|g|g|s|s|s|}
\hline
\multicolumn{5}{|l|}{\textbf{2 -- INSTITUIÇÕES PARCEIRAS}}                                                                                             \\ \hline


{\multirow{2}{*}{Nome completo da Instituição}} &
	{\multirow{2}{*}{Função}} &
	\multicolumn{3}{c|}{Tipo} \\ \cline{3-5} 
\multicolumn{1}{|c|}{} &
	\multicolumn{1}{c|}{} &
	\multicolumn{1}{c|}{Pública} &
	\multicolumn{1}{c|}{Privada} &
	\multicolumn{1}{c|}{ONG} \\ \hline
Casa da Mãe Joana &
	Fornecer abrigo &
	\multicolumn{1}{c|}{\textbf{X}} &
	\multicolumn{1}{c|}{\textbf{}} &
	\multicolumn{1}{c|}{\textbf{}} \\ \hline
Fábrica de Livros &
	Imprimir livros &
	\multicolumn{1}{c|}{\textbf{}} &
	\multicolumn{1}{c|}{\textbf{X}} &
	\multicolumn{1}{c|}{\textbf{}} \\ \hline
Amigos do ambiente &
	vas &
	\multicolumn{1}{l|}{\textbf{}} &
	\multicolumn{1}{l|}{\textbf{}} &
	\textbf{} \\ \hline
		\end{tabularx}
\end{table}

\begin{theo}[APRESENTAÇÃO]
    Este é um modelo para o template de submissão de PIBEX do IFPE \cite{abntex2-wiki-como-customizar}. Aqui tem um exemplo de referência bibliográfica \cite{ibge1993}.
    
    \lipsum[1-3]
\end{theo}

\begin{theo}[CARACTERIZAÇÃO E JUSTIFICATIVA]
    \lipsum[4-6]
\end{theo}

\begin{theo}[OBJETIVOS]
	\lipsum[7]
\end{theo}

\begin{theo}[METODOLOGIA]
	\lipsum[8-12]
\end{theo}

\begin{theo}[RESULTADOS ESPERADOS]
	\lipsum[13-14]
\end{theo}

\begin{theo}[ACOMPANHAMENTO E AVALIAÇÃO]
	\lipsum[15]
\end{theo}

\begin{theo}[CRONOGRAMA]
	\lipsum[16]
\end{theo}

\begin{theo}[RECURSOS E INFRAESTRUTURA EXISTENTES]
	\lipsum[17-18]
\end{theo}

\begin{theo}[RECURSOS E INFRAESTRUTURA NECESSÁRIOS]
	\lipsum[19-20]
\end{theo}

\renewcommand{\bibname}{}

\begin{theo}[REFERÊNCIAS]
\vspace{-1.6cm} 
\bibliography{abntex2-modelo-references}
\end{theo}

\end{document}
